\documentclass[10pt]{article}
\usepackage[margin=1in]{geometry}
\usepackage[T1]{fontenc}
\usepackage{lmodern,microtype,amsmath,amssymb,amsthm}
\theoremstyle{definition}
\newtheorem{assumption}{Assumption}
\begin{document}
\section*{Stochastic Inequality-Constrained Optimization}
We consider minimizing a possibly nonconvex differentiable function $f:E\to\mathbb{R}$ subject to differentiable inequalities $c_i(x)\leq0$, $i=1,\ldots,m$, on a finite-dimensional real inner product space $E$. Gradients are globally Lipschitz and bounded on the feasible set, which satisfies a uniform constraint qualification near active constraints; a given feasible starting point has bounded objective gap. Each oracle call returns stochastic objective gradients, constraint values, and constraint gradients at two points using a shared sample, with independent samples across calls, unbiasedness, bounded variance, and mean-square Lipschitz continuity on feasible points. For $\varepsilon>0$ and $\rho\in(0,1)$, the goal is a feasible output $\bar{x}$ with multipliers $\bar{\lambda}\geq0$ satisfying $\|\nabla f(\bar{x})+\sum_{i=1}^m\bar{\lambda}_i\nabla c_i(\bar{x})\|\leq\varepsilon$ and $\sum_{i=1}^m|\bar{\lambda}_ic_i(\bar{x})|\leq\varepsilon$ with probability at least $1-\rho$.

Write $c=(c_1,\ldots,c_m):E\to\mathbb{R}^m$ and define the feasible set $\mathcal F=\{x\in E:c_i(x)\leq0\text{ for all }i\}$. The problem is
\[
\min_{x\in\mathcal F} f(x).
\]
Let $J(x)=(\nabla c_i(x))_{i=1}^m\in E^m$, with transpose action
\[
J(x)^\top\lambda=\sum_{i=1}^m\lambda_i\nabla c_i(x).
\]
The norm on $E$ is induced by its inner product; $\mathbb{R}^m$ and $E^m$ carry their Euclidean product norms. For triples in $\mathcal V=E\times\mathbb{R}^m\times E^m$, use
\[
\|(a,b,C)\|_{\mathcal V}^2
=\|a\|^2+\|b\|^2+\sum_{i=1}^m\|C_i\|^2.
\]
The parameters satisfy
\[
\begin{aligned}
&m\in\mathbb{N},\quad m\geq1,\quad H\geq1,\quad \sigma,L\geq0,\\
&0<\mu\leq H,\quad \tau>0,\quad \Delta\geq0,\quad
\varepsilon>0,\quad 0<\rho<1.
\end{aligned}
\]
The algorithm receives $E$, the measurable sample space $\Omega$, $m$, $H$, $\sigma$, $L$, $\mu$, $\tau$, $\Delta$, $\varepsilon$, $\rho$, and the starting point $x^0\in E$.

Our assumptions:
\begin{assumption}[Differentiability and gradient bounds]
The functions $f$ and $c_i$, $i=1,\ldots,m$, are Fr\'echet differentiable everywhere. For every $x\in\mathcal F$ and every $i$,
\[
\|\nabla f(x)\|\leq H,
\qquad \|\nabla c_i(x)\|\leq H.
\]
For every $x,y\in E$ and every $i$,
\[
\begin{aligned}
\|\nabla f(x)-\nabla f(y)\|&\leq H\|x-y\|,\\
\|\nabla c_i(x)-\nabla c_i(y)\|&\leq H\|x-y\|.
\end{aligned}
\]
\end{assumption}

\begin{assumption}[Uniform band constraint qualification]
For every $x\in\mathcal F$, there exists $z\in E$ with $\|z\|\leq1$ such that, simultaneously for all $i$,
\[
 c_i(x)\geq-\tau
 \quad\Longrightarrow\quad
 \langle\nabla c_i(x),z\rangle\leq-\mu.
\]
\end{assumption}

\begin{assumption}[Initial feasibility and objective gap]
The given point $x^0$ belongs to $\mathcal F$, and
\[
 f(x^0)-f(x)\leq\Delta,
 \qquad x\in\mathcal F.
\]
\end{assumption}

\begin{assumption}[Stochastic primitive estimates]
There is a probability measure $\nu$ on $\Omega$ and jointly measurable maps
\[
\begin{aligned}
 g &:E\times\Omega\to E,\\
 v &:E\times\Omega\to\mathbb{R}^m,\\
 D &:E\times\Omega\to E^m.
\end{aligned}
\]
Define $S(x,\xi)=(g(x,\xi),v(x,\xi),D(x,\xi))$ and $s(x)=(\nabla f(x),c(x),J(x))$. For every $x\in\mathcal F$, the three components of $S(x,\cdot)$ are Bochner integrable under $\nu$, and
\[
\begin{aligned}
\mathbb{E}_{\xi\sim\nu}[g(x,\xi)]&=\nabla f(x),\\
\mathbb{E}_{\xi\sim\nu}[v(x,\xi)]&=c(x),\\
\mathbb{E}_{\xi\sim\nu}[D(x,\xi)]&=J(x).
\end{aligned}
\]
The squared error is integrable and satisfies
\[
\mathbb{E}_{\xi\sim\nu}
 \bigl[\|S(x,\xi)-s(x)\|_{\mathcal V}^2\bigr]
\leq\sigma^2,
\qquad x\in\mathcal F.
\]
For every $x,y\in\mathcal F$, the squared difference below is integrable, and
\[
\mathbb{E}_{\xi\sim\nu}
 \bigl[\|S(x,\xi)-S(y,\xi)\|_{\mathcal V}^2\bigr]
\leq L^2\|x-y\|^2.
\]
\end{assumption}

At call $k$, a query is any pair $(x,y)\in E\times E$. The oracle returns
\[
\bigl(S(x,\xi_k),S(y,\xi_k)\bigr),
\]
using the same sample for both points and all components. The measurable samples $(\xi_k)_{k\geq0}$ are mutually independent with common distribution $\nu$, and are independent of the algorithm's randomness. One such paired reply counts as one oracle call. The moment conditions above apply to feasible query points.

After the allotted calls, the algorithm must return a measurable pair $(\bar{x},\bar{\lambda})\in E\times\mathbb{R}^m$ satisfying
\[
P\!\left\{
\begin{gathered}
\bar{x}\in\mathcal F,\qquad \bar{\lambda}_i\geq0
\quad(i=1,\ldots,m),\\
\|\nabla f(\bar{x})+J(\bar{x})^\top\bar{\lambda}\|
\leq\varepsilon,\\
\displaystyle\sum_{i=1}^m
 |\bar{\lambda}_i c_i(\bar{x})|\leq\varepsilon
\end{gathered}
\right\}\geq1-\rho,
\]
where $P$ is the joint law of the oracle samples and the algorithm's randomness.
\end{document}
